\section{Integer recourse, flows, and random components}\label{sec:int}

This section treats residual variables that are native integers. The core
search of \cref{sec:rec} applies unchanged, because it uses only the
conditional value function on a fixed residual set. Closure changes. A
patch and a Hessian test are meaningless for integer coordinates; instead,
a few bound-restricted recourse calls certify that one residual integer
label is optimal throughout the retained core hull, and an exact solver for
the small continuous core completes the problem. Under core-only noise
residual ties can persist on every draw, and for integer flows and totally
unimodular systems the label certificate is replaced by one that works with
the whole set of tied optimal labels. The section ends with two results of a
different kind: exact optimization of mixed boxes under strong noise through
random components, and an exact lattice argument for separable integer
problems with a low-rank concave part.

\Cref{tab:int:summary} lists the results. The factor $Q_{\rm ex}$ is defined
in \eqref{eq:rec:Q}, and $c_d$ is the constant of \cref{thm:int:solver}.

\begin{table}[t]
\centering
\caption{Integer recourse and random components. ``Label $+$ algebraic'' is
an integer residual label together with the algebraic output of
\cref{def:model:outputs}\ref{it:model:algebraic} for the core; its size bound
holds on every draw. The function $f_d(k)$ is effective.}
\label{tab:int:summary}
\small
\begin{tabularx}{\textwidth}{@{}lXlll@{}}
\toprule
Result & Noise; residual structure & $b=\log_2M$ & Expected work & Output\\
\midrule
\cref{thm:int:native} & ambient; exact integer oracle on tightened bounds
& $\poly_d(I)$ & $[8^kQ_{\rm ex}+c_d^k]\poly_d(I)$ & label $+$ algebraic\\
\cref{cor:int:native-implicit} & ambient; as above & $\poly_d(I)$
& $8^kQ_{\rm ex}\poly_d(I)$ & label $+$ implicit\\
\cref{thm:int:flow-interior} & core-only; convex integer flow, interior
optimal cores & $\poly_d(I)$ & $[8^kQ_{\rm ex}+c_d^k]\poly_d(I)$
& label $+$ algebraic\\
\cref{thm:int:flow-boundary} & core-only; convex integer flow
& $f_d(k)\poly_d(I)$ & $f_d(k)Q_{\rm ex}\poly_d(I)$ & label $+$ algebraic\\
\cref{cor:int:bilinear} & core-only; flow, bilinear coupling & $\poly_d(I)$
& $[8^kQ_{\rm ex}+c_d^k]\poly_d(I)$ & label $+$ algebraic\\
\cref{thm:int:tu} & core-only; totally unimodular & as for flows
& as for flows & label $+$ algebraic\\
\cref{thm:int:strong-field} & ambient, strong; any mixed box & any
& see \eqref{eq:int:sf-work} & component\\
\cref{thm:int:lowrank} & aligned; separable integer box & $\poly(I)$
& $8^kH_{\rm rat}\poly(I)$ & rational, no fallback\\
\bottomrule
\end{tabularx}
\end{table}

\subsection{An exact solver for a small continuous core}\label{sec:int:solver}

All label certificates end with the same task: minimize an explicit
polynomial exactly over a box of small dimension. A general exact method
returns separate defining polynomials for the coordinates; pairing their
roots costs a product of degrees, and it does not by itself say which
coordinate roots belong to the same optimizer. The following solver returns
all coordinates as polynomials in one common real algebraic number, with a
base that is exponential only in the dimension.

\begin{theorem}[Exact box solver through joint critical limits]\label{thm:int:solver}
Fix $d\ge1$, and let $D$ be the least even integer greater than $d$. There
are a constant $c_d$, depending only on $d$, and a deterministic algorithm
with the following properties. The input is an explicit rational polynomial
$f$ of degree at most $d$ in $k$ variables and a nonempty closed rational box
$B\subseteq\R^k$; let $H$ be its length. The algorithm returns a global
minimizer $x^*$ of $f$ on $B$ and the value $f(x^*)$ in the algebraic format
of \cref{def:model:outputs}: the coordinates whose interval is a point are
rational, and the others are $x^*_i=r_i(\theta)$ for polynomials
$r_i\in\Q[T]$ and one real root $\theta$ of a squarefree integer polynomial
$P$ of degree at most $(D-1)^k$, given with a rational isolating interval;
the value is given by its own defining polynomial and isolating interval.
The bit work and the output length are at most $c_d^k\poly_d(H)$. For every
$q\ge0$, a rational point within distance $2^{-q}$ of $x^*$, a feasible
rational point with certified objective gap at most $2^{-q}$, and a rational
interval of width at most $2^{-q}$ containing $f(x^*)$ are computed within
$c_d^k\poly_d(H+q)$. No genericity is assumed: ties, singular critical points,
and positive-dimensional sets of minimizers are allowed.
\end{theorem}

For each face of the box, the solver deforms the restricted objective by
$\varepsilon\sum_ix_i^D$ with a symbolic $\varepsilon$. The deformed stationary
equations have leading monomials $x_i^{D-1}$, so they form a Gr\"obner basis
for every value of $\varepsilon$, and the quotient algebra has the fixed
dimension $(D-1)^r$ on a face with $r$ free coordinates. The leading
coefficient, in $1/\varepsilon$, of the characteristic polynomial of a linear
form $\lambda^{\mathsf T}x$ in the quotient vanishes exactly at the
projections of the finite limits of critical points as $\varepsilon\to0$, and
its derivatives with respect to the coefficients of $\lambda$ recover all
coordinates of each limit as polynomials in its projection. A deterministic
list of $O(r(D-1)^{2r})$ moment-curve forms contains one that works; the
others can only add feasible candidates, which are harmless. Minimizers of
the deformed problems converge to a global minimizer of $f$, so the candidate
of least value is optimal. The proof is in \cref{app:int:solver}; it uses
finite-dimensional quotient algebras, simultaneous triangularization, Puiseux
series, and exact univariate algorithms, and \cref{lem:int:budget} gives the
explicit choice $c_d=2^{10000D}$ for one elementary implementation. That
constant is a proved bound used to define noise laws below, not a practical
estimate. The ingredients are classical. Recovering all coordinates from one
univariate polynomial and derivatives of a characteristic polynomial is the
mechanism of the rational univariate representation \citep{rouillier1999-solving-zero-dimensional-systems-through},
and deformations of stationary systems are standard in critical-point
methods \citep{basu2006-algorithms-in-real-algebraic-geometry}. We give a self-contained proof for boxes,
including singular and positive-dimensional critical sets, because the
recourse theorems need its exact contract and its constant-base bound; we do
not claim a new general algebraic complexity result.

\begin{corollary}[Mixed boxes and exact comparisons]\label{cor:int:mixed-solver}
For a mixed box with $k$ continuous coordinates and native integer
coordinates with $N_1,\ldots,N_t$ labels, enumerating the labels and applying
\cref{thm:int:solver} to each gives an exact global minimizer and value in
$c_d^k\prod_iN_i\poly_d(H)$ bit operations. Values of two algebraic outputs
of \cref{thm:int:solver} can be compared exactly, including equality, within
$c_d^k\poly_d(H)$; the comparison uses the two univariate value polynomials
and forms no field containing both.
\end{corollary}

\subsection{Native integer recourse}\label{sec:int:native}

\begin{definition}[Native integer recourse]\label{def:int:native}
A \emph{native integer recourse instance} is a core--recourse instance
(\cref{def:rec:instance}) whose residual set is
\[
 Y=\{z\in\Z^{n_R}:\ Az\le b,\ \ell\le z\le u\},
\]
with a rational matrix $A$, a rational vector $b$, and integer bounds
encoded in binary; the curvature bound \eqref{eq:rec:L} is required on
$[0,1]^k\times\prod_i[\ell_i,u_i]$. Put $N_Z=\prod_i(u_i-\ell_i+1)$ and
$S=k+\sum_i(u_i-\ell_i)$. A \emph{box-stable exact integer oracle} receives
a rational core point $v$, rational linear coefficients, and integer bounds
$\ell\le\ell'\le u'\le u$, and returns the exact optimal value of
$\min\{F_\gamma(v,z):z\in Y,\ \ell'\le z\le u'\}$ and an attaining integer
point, or reports infeasibility, in bit work polynomial in the query length
with an exponent independent of $k$.
\end{definition}

The core changes costs, never feasibility: supplies, right-hand sides, and
bounds do not depend on $v$. Exact values are rational because the core point
is rational and residual points are integral. The oracle must handle
arbitrary tightened coordinate bounds, which is how competing labels are
excluded.

\begin{theorem}[Native integer recourse]\label{thm:int:native}
Fix $d$. Let a native integer recourse instance with $k\ge1$ be given,
together with a box-stable exact integer oracle. There is a power of two $M$,
computed from the base instance with $\log_2M\le\poly_d(I)$, such that for the
ambient perturbation model with marginal law $U_{\sigma,M}$ the following
hold.
\begin{enumerate}[label=(\roman*)]
\item On every draw the algorithm returns a global minimizer $(v^*,z^*)$ of
$F_\gamma$ on $X=[0,1]^k\times Y$ and the optimal value: the label $z^*$
exactly, and $v^*$ and the value in the algebraic format of
\cref{thm:int:solver}. Its length is at most $c_d^k\poly_d(I)$ on every
draw, and a $q$-bit refinement costs $c_d^k\poly_d(I+q)$.
\item The expected bit work and the expected size of the global proof record
are at most $[8^kQ_{\rm ex}+c_d^k]\poly_d(I)$.
\item If $F_0$ is quadratic, the output is rational on every draw and the
expected work is $[8^kQ_{\rm ex}+3^k]\poly(I)$.
\end{enumerate}
No growth, label-gap, uniqueness, or residual convexity premise is used.
\end{theorem}

The certificate is the following. At the incumbent corner $c_j$ with label
$z_j$, the $2n_R$ restrictions $z_i\le z_{j,i}-1$ and $z_i\ge z_{j,i}+1$,
each together with the original constraints, have union exactly
$Y\setminus\{z_j\}$, however the residual coordinates are coupled. Let
$V_{\rm other}(c_j)$ be the least of their exact values. If
\begin{equation}\label{eq:int:label}
 V_{\rm other}(c_j)-V(c_j)>2G\,w_\infty(\mathcal H_j),
\end{equation}
then $z_j$ is the unique optimal label at every core point of the retained
hull $\mathcal H_j$ (\cref{lem:rec:value}(c)), so every global minimizer lies
in $\mathcal H_j\times\{z_j\}$. The whole slice $[0,1]^k\times\{z_j\}$ is
feasible and contains a global minimizer, so \cref{thm:int:solver} applied to
$F_\gamma(\cdot,z_j)$ on $[0,1]^k$ finishes the problem, whether the core
minimizers are unique, tied, or form a continuum. Feasible values for the
restricted problems would not suffice; \eqref{eq:int:label} needs their exact
optimal values. Distinct integer labels are at unit distance, so point growth
forces \eqref{eq:int:label} at a base-computed level; no label gap is supplied
or charged separately. On the remaining draws the algorithm enumerates the
$N_Z$ labels on the same draw, solves each core slice, and keeps the best by
exact comparisons (\cref{cor:int:mixed-solver}). It returns only the winning
label and its core representation, so the output bound does not depend on
how many labels were examined. The proof is in \cref{app:int:native}.

\begin{corollary}[Implicit core output]\label{cor:int:native-implicit}
Under the hypotheses of \cref{thm:int:native}, an algorithm with expected bit
work $8^kQ_{\rm ex}\poly_d(I)$ returns on ordinary draws the exact label $z^*$
and an implicit patch output for the core slice $F_\gamma(\cdot,z^*)$, and on
the remaining draws the output of \cref{thm:int:native}. A $q$-bit evaluation
of the ordinary output costs $\poly_d(I+q)$, and its expected cost over all
draws is $\poly_d(I+q)$.
\end{corollary}

The variant replaces the exact core solve by the sign and Hessian tests of
\cref{lem:rec:closure} on the core coordinates of the certified slice; it pays
for an additional active-gradient event instead of the term $c_d^k$.

\paragraph{Separable convex recourse over totally unimodular systems.}
The main source of exact integer oracles in this paper is the following
classical result.

\begin{proposition}[Separable convex integer oracles]\label{prop:int:hs}
Let $Y=\{z\in\Z^{n_R}:Az=b,\ \ell\le z\le u\}$, where $A$ is an integer
totally unimodular matrix, $b$ is integral, and the bounds are integral and
finite. Let $F_0(v,z)=\phi(v)+\sum_if_i(v,z_i)$ with explicit rational
polynomials of fixed degree, and suppose $t\mapsto f_i(v,t)$ is convex on
$[\ell_i,u_i]$ for every $v\in[0,1]^k$ (promised, or certified). Then a
box-stable exact integer oracle exists for every rational linear perturbation,
with bit work polynomial in the query length and logarithmic dependence on the
numerical bounds. This includes integer flows on a directed network with
integral supplies and arc bounds; inequality systems $Cz\le b$ with $C$ totally
unimodular reduce to this form (\cref{cor:int:tu-ineq}).
\end{proposition}

The oracle is the proximity-scaling algorithm of
\citet[Theorem~4.3]{hochbaum1990-convex-separable-optimization-is-not}, which solves separable convex
integer minimization over totally unimodular systems with logarithmically
many linear programs in the numerical bounds and function evaluations at
prescribed points. At a rational core point all evaluations are rational of
polynomial length; tightening coordinate bounds and adding linear terms keep
the class; and convex tangent extensions beyond the native intervals supply
evaluations if the algorithm requests them
(\cref{app:int:native}). The convexity premise concerns each scalar cost on its
whole real interval. For networks, optimality also has a short certificate.

\begin{lemma}[Marginal potentials]\label{lem:int:potentials}
Let $Y$ be the integer flows of a directed network with fixed integral
supplies and arc bounds, and let each $t\mapsto f_a(t)$ be convex. A feasible
$z$ minimizes $\sum_af_a(z_a)$ over $Y$ if and only if there are node
potentials $\pi$ for which every arc of the residual network of $z$ has
nonnegative reduced cost, where a forward arc $a=(p,p')$ with $z_a<u_a$ costs
$f_a(z_a+1)-f_a(z_a)$ and a backward arc with $z_a>\ell_a$ costs
$f_a(z_a-1)-f_a(z_a)$, and the reduced cost adds $\pi_p-\pi_{p'}$ along the
residual arc from $p$ to $p'$. Shortest-path distances from an added zero-cost
source supply such potentials of polynomial bit length.
\end{lemma}

This is a certificate, not an algorithm: repeated unit augmentation is not
claimed to run in polynomial time. With \cref{prop:int:hs},
\cref{thm:int:native} applies to separable convex integer flows and totally
unimodular systems whose costs depend arbitrarily on the core. For instance,
$F_0(v,z)=\phi(v)+\sum_a\psi_a(z_a)+v^{\mathsf T}Wz$ with convex $\psi_a$ allows
arbitrary core-to-arc coupling $W$; the curvature bound $L$ is that of $\phi$,
and $W$ enters only through its bit length. Convexity of a fixed-degree
univariate $\psi_a$ on its interval can be verified by univariate sign
determination.

\subsection{Core-only noise for integer flows}\label{sec:int:flows}

Under core-only noise the residual costs are not perturbed, and many optimal
labels can persist on every draw. Full point growth then fails, and the
competing-label certificate cannot pass when two labels are optimal at the
optimal core. Throughout this subsection $Y$ is the set of integer flows of a
directed network with $s$ nodes and $n_R$ arcs, integral supplies, and
integral arc bounds $[\ell_a,u_a]$ in binary, and
\begin{equation}\label{eq:int:flow-model}
 F_\gamma(v,z)=\phi(v)+\sum_af_a(v,z_a)+\gamma^{\mathsf T}v,
 \qquad v\in[0,1]^k,\ z\in Y,
\end{equation}
with explicit rational polynomials of degree at most $d\ge1$, each
$f_a(v,\cdot)$ convex on $[\ell_a,u_a]$ for every $v$ (promised or certified),
and a curvature bound $L$ as in \eqref{eq:rec:L}. Only the $k$ core
coefficients are perturbed. By \cref{prop:model:regret} the sampled optimizer
is within $k\sigma$ of $\min_XF_0$ in original objective value, whatever the
size of the network.

A natural certificate selects one flow and proves it optimal on the retained
hull by polynomial potentials. Boundary core optima defeat it:
\cref{prop:lim:flow} gives two parallel arcs and a two-dimensional core on
which, with probability $\frac1{16}$ for every grid size, the optimal core
is a vertex with projected growth constant at least one, both flows are
optimal there, and no single flow is optimal on any box at that vertex,
while the core search retains such a box at every level. With $m$ stages,
an algorithm that insists on a single-flow certificate and enumerates
labels when it fails does work at least $2^m/16$ in expectation. The
obstruction concerns that certificate; the instance itself is easy. Two
remedies follow: a premise that all optimal cores are interior, and a
certificate that uses all tied flows at once.

\begin{theorem}[Core-only flow recourse with interior optimal cores]\label{thm:int:flow-interior}
Fix $d$, and assume, in addition to the model \eqref{eq:int:flow-model} with
$k\ge1$, that for every $\gamma\in[-\sigma,\sigma]^k$ every optimal core point
of $F_\gamma$ lies in $(0,1)^k$ (promised, or certified, for instance by
$\partial_iF_0\le-\sigma-\tau'$ on $\{v_i=0\}$ and $\partial_iF_0\ge\sigma+\tau'$
on $\{v_i=1\}$ for all feasible flows and some rational $\tau'>0$). There is a
power of two $M$, computed from the base instance with $\log_2M\le\poly_d(I)$,
such that for core-only noise with marginal law $U_{\sigma,M}$ the algorithm
returns on every draw an optimal integer flow and an algebraic output of one
optimal core and the optimal value, of length $c_d^k\poly_d(I)$, with expected
bit work $[8^kQ_{\rm ex}+c_d^k]\poly_d(I)$. For quadratic objectives the output
is rational.
\end{theorem}

The interiority premise is used only in the running-time analysis;
correctness holds on every draw without it, and its sufficient certificate
enters no numerical factor. At the incumbent corner the algorithm computes an
optimal flow, a tight shortest-path tree of its residual network, and the
corresponding polynomial potentials, and it tests exactly whether every
residual reduced cost is nonnegative on the whole retained hull; identically
zero reduced costs pass, so persistent ties are allowed. If the test passes,
the flow is optimal on the hull, and the exact solver finishes its whole core
slice. The analysis uses a finite family of reduced-cost polynomials over all
labels and trees, fixed before sampling: an optimal core in the interior is a
stationary point of every attaining fixed-flow slice, so the noise equals a
negative gradient there, and the event that the optimal core is near a zero
of a chart polynomial has small probability by an algebraic tube estimate.

\paragraph{All tied flows at once.}
At a core point $c$ on a face $\mathcal A$ of the core cube, let $z^0$ be an
optimal flow and $\pi$ potentials certifying it (\cref{lem:int:potentials}).
The adjusted costs $g_a(c,t)=f_a(c,t)+(\pi_p-\pi_{p'})t$ of the arcs $a=(p,p')$
are convex, and their integer minimizers form intervals $I_a\ni z^0_a$, found
by binary search on monotone first differences. The set
$Y_0=\{z\in Y:z_a\in I_a\text{ for all }a\}$ is exactly the set of optimal flows
at $c$: the potential terms sum to a constant over $Y$, and the objective gap is
a sum of nonnegative scalar gaps. Minimizing an inward core derivative over
$Y_0$ is again a convex flow problem, because on an interval with at least
three integers the adjusted cost is constant, so the original arc cost is
affine there and its inward derivative is convex; two-point intervals use
linear interpolation, and singleton intervals are substituted. Finally, every
feasible flow is close to $Y_0$.

\begin{lemma}[Proximity to the optimal set]\label{lem:int:proximity}
For every $z\in Y$ there is $\bar z\in Y_0$ with
$\norm{z-\bar z}_1\le n_R\sum_a\dist(z_a,I_a)$.
\end{lemma}

The deterministic certificate uses these facts on a rational core box
$\mathcal D$ whose intervals in the normal coordinates of $\mathcal A$ contain
the corresponding bounds, with projection $\mathcal D_{\mathcal A}$ onto the
face and $c\in\mathcal D_{\mathcal A}$. Let $\varsigma_i=+1$ for a normal
coordinate fixed at $0$ and $\varsigma_i=-1$ for one fixed at $1$, let $T_1$ and
$T_\infty$ be the sum and the maximum of the normal widths of $\mathcal D$, let
$R_{\mathcal A}$ bound $\norm{w-c}$ on $\mathcal D_{\mathcal A}$, and let
$H\ge1$ and $K\ge0$ be rational bounds on $\norm{\nabla^2_vF_0}$ and on
$\abs{\partial_{v_i}\partial_{z_a}f_a}$ on the real bounding domain.

\begin{theorem}[An optimal-face certificate]\label{thm:int:face}
Suppose that, with the potentials taken as polynomials of the core obtained
from a tight shortest-path tree at $c$:
\begin{enumerate}[label=(\alph*)]
\item for every arc, $g_a(w,t)=g_a(w,z^0_a)$ for all $w\in\mathcal D_{\mathcal A}$
and all integers $t\in I_a$;
\item every first adjusted marginal outside $I_a$ that stays in the native
interval is at least $\mu_{\rm out}=n_RKT_1$ on $\mathcal D_{\mathcal A}$; and
\item if $\mathcal A$ has normal coordinates, $\beta=\min_i\beta_i$ satisfies
$\beta-HR_{\mathcal A}-\frac H2T_\infty>0$, where
$\beta_i=\min_{z\in Y_0}\varsigma_i\partial_{v_i}F_\gamma(c,z)$.
\end{enumerate}
Then every flow of $Y_0$ is optimal at every point of $\mathcal D_{\mathcal A}$,
and every point of $\mathcal D\times Y$ off the face is strictly worse than some
feasible point on it. If $\mathcal D$ contains the core of every global
minimizer, the slice $[0,1]^k\times\{z^0\}$ contains a global minimizer.
Conditions (a)--(c) are decided exactly, (a) by at most $d+1$ polynomial
identities per arc, (b) by exact polynomial minimization on a rational box
(\cref{thm:int:solver}), and (c) by one convex flow problem per normal
coordinate.
\end{theorem}

The certificate is sound on every draw and needs neither uniqueness of the
optimal flow nor a noise assumption. Both (b) and (c) are needed: (c) uses the
minimum over \emph{all} optimal flows, since a single tied flow can have the
wrong inward sign, and without (b) a flow that loses at $c$ can become optimal
off the face. In the instance of \cref{prop:lim:flow} the certificate with the origin
face closes every box of maximum width less than $2\min_i\gamma_i$.

\begin{theorem}[Core-only flow recourse at arbitrary core faces]\label{thm:int:flow-boundary}
Fix $d$. For the model \eqref{eq:int:flow-model} with $k\ge1$ there are an
effective function $f_d$ and a power of two $M$, computed from the base
instance with
\[
 \log_2M\le f_d(k)\poly_d(I),
\]
such that for core-only noise with marginal law $U_{\sigma,M}$ the algorithm
returns on every draw an optimal integer flow and an algebraic output of one
optimal core and the optimal value, of length $f_d(k)\poly_d(I)$, with
$q$-bit refinement in $f_d(k)\poly_d(I+q)$, and with expected bit work
$f_d(k)Q_{\rm ex}\poly_d(I)$. No interiority, growth, or uniqueness premise is
used, and arbitrarily many optimal flows may persist.
\end{theorem}

At each level the algorithm tries the certificate of \cref{thm:int:face} on
every face of the core cube compatible with the retained hull, at the midpoint
of the projected hull, and solves the whole core slice of $z^0$ when one passes.
The analysis needs three probability bounds: projected growth, distance of the
optimal core from the zeros of a base-only family of chart polynomials (tube
estimate), and a margin $\tau$ for the minimal inward derivative over all
optimal flows. The last is obtained without conditioning on a random label: on
each fixed face, the optimal core and its whole optimal flow set do not depend
on the normal noise coefficients. A tube controls the distance from a zero set,
not the value of a polynomial away from it. Converting distance into a value
margin uses fixed-block quantifier elimination with its coefficient-height
bound, and costs a precision of $f_d(k)\poly_d(I)$ bits. This is why the
sampling precision in \cref{thm:int:flow-boundary} depends on $k$.

\begin{corollary}[Bilinear coupling]\label{cor:int:bilinear}
If $F_\gamma(v,z)=\phi(v)+\sum_a\psi_a(z_a)+v^{\mathsf T}Wz+\gamma^{\mathsf T}v$
with a rational matrix $W$, fixed degree $d\ge2$, convex $\psi_a$ on
$[\ell_a,u_a]$, and $\partial_{ii}\phi\le L$ on $[0,1]^k$, then the conclusions
of \cref{thm:int:flow-boundary} hold with $\log_2M\le\poly_d(I)$, output length
$c_d^k\poly_d(I)$, refinement $c_d^k\poly_d(I+q)$, and expected bit work
$[8^kQ_{\rm ex}+c_d^k]\poly_d(I)$. For quadratic $\phi$ the output is rational.
\end{corollary}

All chart polynomials are then affine in the core, and an affine polynomial
satisfies $\abs{p(x)}\ge b_{\min}\dist(x,Z(p)\cap[0,1]^q)$, with $b_{\min}$ its
smallest nonzero coefficient, which restores a polynomial-length value margin.
The curvature in the numerical factor is that of $\phi$ alone. This would fail
for nonlinear coupling: a term $v_i^2z_a^2$ makes the core curvature grow
quadratically with the capacity of arc $a$.

\subsection{Totally unimodular systems}\label{sec:int:tu}

\begin{theorem}[Totally unimodular recourse]\label{thm:int:tu}
Let $Y=\{z\in\Z^{n_R}:Az=b,\ \ell\le z\le u\}$ be nonempty, with $A$ an integer
totally unimodular matrix (promised), $b$ integral, and finite integral bounds,
and let $F_\gamma(v,z)=\phi(v)+\sum_if_i(v,z_i)+\gamma^{\mathsf T}v$ with
$f_i(v,\cdot)$ convex on $[\ell_i,u_i]$ for every $v$. Then the conclusions of
\cref{thm:int:flow-boundary} hold. If moreover
$F_\gamma=\phi(v)+\sum_i\psi_i(z_i)+v^{\mathsf T}Wz+\gamma^{\mathsf T}v$ as in
\cref{cor:int:bilinear}, its conclusions hold.
\end{theorem}

Two ingredients replace network structure. Shortest-path potentials become a
dual vector satisfying at most $2n_R$ adjacent-slope inequalities, which exist
because the piecewise-linear interpolation of the separable costs has an
integral optimum over the real totally unimodular polytope; a dual vertex in a
pre-draw box has inverse-basis entries in $\{0,\pm1\}$, which bounds the chart
coefficients. Directed cycles become conformal circuits of $\ker A$ with entries
in $\{0,\pm1\}$, which gives \cref{lem:int:proximity} with the same factor
$n_R$. Both ingredients use total unimodularity; neither holds for general
integer matrices.

\begin{corollary}[Bounded totally unimodular inequalities]\label{cor:int:tu-ineq}
The conclusions of \cref{thm:int:tu} hold for
$Y=\{z\in\Z^{n_R}:Cz\le b,\ \ell\le z\le u\}$ with $C$ totally unimodular, $b$
integral, and finite integral bounds.
\end{corollary}

Add the slack $s_i=b_i-C_iz$ with the explicit bounds
$0\le s_i\le b_i-\min_{\ell\le z\le u}C_iz$; the matrix $[C\ I]$ is totally
unimodular, projection is a bijection onto $Y$, and slacks have zero cost, so
the curvature bound and all objective values are unchanged. A negative slack
bound proves infeasibility.

\subsection{Strong noise and random components}\label{sec:int:strong-field}

The previous results need small noise and a small core. The following result
needs neither a core nor a decomposition, but it needs noise that dominates
the variation of every coordinate derivative. Let $X=\prod_iX_i$ be a mixed
box (\cref{sec:model:input}) with $n$ coordinates after preprocessing, and let
$F_0$ be an explicit rational polynomial of degree at most $d$. The
\emph{primal graph} joins two coordinates when they occur together in a
nonzero monomial, so each monomial support is a clique; let
$\Delta_+=\max\{1,\text{maximum degree}\}$. Before sampling, compute rational
enclosures $\partial_iF_0\in[\underline\partial_i,\overline\partial_i]$ on the
continuous hull, by multiplying exact interval ranges of univariate powers
over each derivative monomial. Put $R_i=\overline\partial_i-\underline\partial_i$,
\[
 a_i=\begin{cases}c_d,&i\text{ continuous},\\u_i-\ell_i+1,&i\text{ native integer},
 \end{cases}\qquad
 q_i=U_{\sigma,M}\bigl([-\overline\partial_i,-\underline\partial_i]\bigr),\qquad
 \beta=\max_ia_iq_i,
\]
where $c_d$ is the constant of \cref{thm:int:solver}; for quadratic
objectives one may take $a_i=3$ for continuous coordinates. Each $q_i$ is an
exactly computable rational number.

\begin{theorem}[Strong fields and random components]\label{thm:int:strong-field}
Fix $d$, and let $M\ge2$ be any power of two. Under the ambient perturbation
model with marginal law $U_{\sigma,M}$, a deterministic algorithm returns on
every draw an exact global minimizer of $F_\gamma$ on $X$ in the component
format of \cref{def:model:outputs}\ref{it:model:component}: rational values
for the pinned coordinates, and for each remaining component its integer labels
and one common-root algebraic representation of its continuous coordinates and
value. If $4\Delta_+\beta<1$, its expected bit work is at most
\begin{equation}\label{eq:int:sf-work}
 \poly_d(I+b)\Bigl[1+\frac{\sum_ia_iq_i}{1-4\Delta_+\beta}\Bigr],\qquad b=\log_2M,
\end{equation}
and $q$-bit evaluations of the optimal value and of a feasible rational point
with certified total gap $2^{-q}$ cost the same with $\poly_d(I+b+q)$. For
quadratic objectives the output is an exact rational optimizer and value, with
$a_i=3$ for continuous coordinates. A sufficient condition for a bounded
denominator is
\begin{equation}\label{eq:int:sf-regime}
 \sigma\ge8\Delta_+\max_i(a_iR_i),\qquad M\ge16\Delta_+\max_ia_i,
\end{equation}
which gives $a_iq_i\le1/(8\Delta_+)$ and $4\Delta_+\beta\le\frac12$.
\end{theorem}

Outside the closed interval in the definition of $q_i$, the sampled
derivative $\partial_iF_\gamma$ has a strict sign throughout the continuous
hull, so every global minimizer has $x_i$ at the corresponding bound; this
holds for integer coordinates as well, by monotonicity between labels. These
\emph{bad} events depend on one coefficient each, so they are independent
Bernoulli events with the exact probabilities $q_i$; recomputing the
enclosures after pinning would destroy this independence and is not needed.
After pinning, every surviving monomial lies in one connected component of the
graph induced by the bad coordinates, so the objective splits into independent
component problems, each solved exactly by \cref{cor:int:mixed-solver} or, for
quadratics, by face enumeration. At most $(4\Delta_+)^{t-1}$ connected sets of
size $t$ contain a given vertex, and a subcritical branching bound gives
\eqref{eq:int:sf-work}. The proof is in \cref{app:int:strong-field}.

The algorithm is exact on every draw, also when $4\Delta_+\beta\ge1$, but the
theorem then gives no useful bound. Condition \eqref{eq:int:sf-regime} is a
genuinely strong-noise regime. It scales with the full derivative variation,
with the graph degree, and, for integer coordinates, with the numerical number
of labels; for continuous coordinates of nonquadratic objectives it scales
with the large effective constant $c_d$. The total value is reported as a
rational constant plus a sum of component values. Expanding that sum into one
defining polynomial can have degree exponential in the number of components,
and exact sign or equality tests for sums of unrelated algebraic numbers are
not claimed; rational enclosures of the total are obtained by refining each
component. The regret bound of \cref{prop:model:regret} is
$\sigma\sum_iw_i$, which is large in this regime.

\subsection{Separable integer problems with a low-rank concave part}\label{sec:int:lowrank}

The last result uses aligned noise (\cref{def:model:perturbation}) and an exact
argument of a different kind: the original objective values of the sampled
problem lie on a lattice whose spacing is controlled by the same grid that
controls the expected count. Let $X=\prod_{j=1}^n\{L_j,\ldots,U_j\}$ with
integer bounds in binary, and
\begin{equation}\label{eq:int:lowrank-model}
 F(x)=\sum_{j=1}^ng_j(x_j)-\frac\alpha2\norm{Tx}^2,
\end{equation}
where each $g_j$ is a rational polynomial of degree at most four, convex on
$[L_j,U_j]$ (verified by minimizing its quadratic second derivative), $\alpha>0$
is rational, and $T\in\Q^{k\times n}$ is supplied. For each row choose a
rational $\sigma_i>0$, and let $\ell_i=\min_X(Tx)_i$, $u_i=\max_X(Tx)_i$,
$w_i=u_i-\ell_i+2\sigma_i/\alpha$, $s=\max_iw_i$. Let $Q_d$ be the product of
the denominators of all base rational data, put $D_0=2Q_d^3$, and let $M$ be the
least power of two with $M\ge\max\{2,k\alpha s^2D_0\}$. The sampled objective
is $F_\xi(x)=F(x)+\xi^{\mathsf T}Tx$, where $\xi_1,\ldots,\xi_k$ are independent
and $\xi_i$ has law $U_{\sigma_i,M}$.

\begin{theorem}[Exact low-rank separable integer optimization]\label{thm:int:lowrank}
For every draw, a deterministic algorithm returns an exact integer minimizer
of $F_\xi$ on $X$ and the exact rational optimal value. Its expected bit work is
at most
\[
 8^kH_{\rm rat}\poly(I),\qquad
 H_{\rm rat}=\prod_{i=1}^k\Bigl[3+\frac{(1+2k)\alpha w_i}{2\sigma_i}\Bigr],
\]
with an absolute polynomial exponent. No growth, uniqueness, or separation
premise is used, there is no fallback, and the number of integer coordinates
and the numerical lengths of their intervals do not enter $H_{\rm rat}$ except
through the row ranges $u_i-\ell_i$.
\end{theorem}

Completing the square gives the auxiliary value
$W(a)=\frac\alpha2\norm a^2+\min_{x\in X}[G(x)-\alpha a^{\mathsf T}Tx]$ on the
fixed box $\prod_i[\ell_i-\sigma_i/\alpha,u_i+\sigma_i/\alpha]$, with
$G=\sum_jg_j$. It has upper coordinate curvature $\alpha$, its inner problem
separates into scalar convex integer problems solved by binary search on
monotone differences, and $\min_aW(a)+\xi^{\mathsf T}a$ equals $\min_XF_\xi$ up to
the known constant $\norm\xi^2/(2\alpha)$. The core search of
\cref{sec:rec:value} on the $k$ auxiliary coordinates gives an integer witness
whose original gap is below $E_J<1/(D_0(M-1))$ at the predetermined level
$J=\log_2M$, while every original value $F_\xi(x)$ lies in
$\frac1{D_0(M-1)}\Z$: one common grid denominator suffices, because the noise
enters the original objective linearly. The auxiliary values may have
denominators involving $(M-1)^2$, but no lattice property of them is used. The
proof is in \cref{app:int:lowrank}.

The numerical ratios $\alpha w_i/\sigma_i$ are explicit; binary-encoded bounds
alone do not make them small. For binary variables, and more generally for
explicitly listed labels, deterministic fixed-rank algorithms already exist: on
$\{0,1\}^n$ the objective is concave on the zonotope image of
$x\mapsto(Tx,c^{\mathsf T}x)$, so a vertex is optimal, and the zonotope has
$O(n^k)$ vertices; with explicit labels the planar lifted label polytopes give
a Minkowski-sum enumeration polynomial in the total label count. The theorem
does not improve these cases; it allows binary-encoded intervals without listing
their labels.

\subsection{Scope}\label{sec:int:scope}

The residual feasible set must not depend on the core, and the exact integer
oracle must handle all tightened coordinate bounds. The flow and totally
unimodular results need separable costs that are convex in each residual
coordinate for every core point; they do not cover nonseparable integer
convexity, arbitrary integer matrices, or core-dependent supplies and
right-hand sides. Under core-only noise, single-flow certificates fail at
boundary cores (\cref{prop:lim:flow}); the face certificate needs a
parameter-dependent sampling precision for nonlinear coupling, which bilinear
coupling avoids. The strong-field theorem is a strong-noise statement, and the
low-rank theorem uses aligned rather than ambient noise. All results optimize
the sampled objective exactly; none optimizes $F_0$. The integer oracles,
fixed-block elimination, tube estimates, and connected-set counting are
established tools; what is proved here is their composition with one
base-selected finite law, the label and face certificates, and the bit and
output bounds, including exact algebraic output that shares one root for all
coordinates of the optimizer.
